\section{Current Limitations}
\label{sec:limitations}

Despite substantial progress in applying generative models to rare disease research, significant challenges remain before these approaches achieve routine clinical deployment. Validation frameworks lack standardization, privacy guarantees remain incomplete, computational requirements constrain accessibility, and the translation from research prototypes to regulatory-approved diagnostic tools faces substantial barriers. This section examines these limitations systematically to identify priorities for future development.

\subsection{Validation and Evaluation Challenges}

The absence of standardized validation protocols represents perhaps the most fundamental limitation constraining generative model deployment. Finetti et al. identified that over 40\% of reviewed publications provided insufficient quantitative detail regarding dataset characteristics, while fewer than 33\% reported both pre- and post-augmentation metrics \cite{finetti2025data}. This inconsistent reporting impedes meta-analysis, reproducibility assessment, and comparative evaluation across different generative approaches. Without consensus evaluation frameworks, researchers cannot reliably determine whether observed performance differences reflect genuine methodological advances or artifacts of dataset selection, hyperparameter choices, or evaluation protocols.

Current validation practices predominantly rely on downstream task performance—measuring classification accuracy, segmentation quality, or prediction metrics on models trained with synthetic data. While pragmatic, this approach provides only indirect evidence of synthetic data quality. Strong performance may result from the model memorizing training patterns rather than learning generalizable representations. Conversely, weak performance could indicate limitations in the downstream model architecture rather than synthetic data quality. Finetti et al. noted that biological plausibility assessment remains rare, with most studies evaluating synthetic samples solely through quantitative metrics rather than domain expert review \cite{finetti2025data}.

Visual inspection by clinical experts, employed successfully by Frid-Adar et al. for liver lesion validation \cite{fridadar2018gan}, does not scale to large synthetic datasets and introduces subjectivity that complicates systematic evaluation. Automated quality assessment metrics specific to biomedical domains remain underdeveloped. General-purpose image quality metrics such as Structural Similarity Index (SSIM) or Fréchet Inception Distance (FID) may not capture medically-relevant features that determine diagnostic utility. For electronic health records, statistical tests comparing marginal distributions fail to detect higher-order dependencies or clinically impossible code combinations.

The temporal dimension adds further complexity. For longitudinal health records or disease progression modeling, validation must assess not only cross-sectional realism but also the plausibility of temporal trajectories. A synthetic patient record might contain individually realistic diagnosis codes yet exhibit an impossible progression sequence—for example, recording diabetes complications before the initial diabetes diagnosis. Zhong et al. addressed this through conditioning mechanisms that respect temporal ordering \cite{zhong2024meddiffusion}, but systematic evaluation frameworks for longitudinal plausibility remain absent.

\subsection{Privacy and Security Vulnerabilities}

Researchers have promoted synthetic data generation as a privacy-preserving alternative to traditional de-identification approaches, yet substantial vulnerabilities persist. Mendes et al. documented that generative models remain susceptible to membership inference attacks, where adversaries determine whether the training set included a specific individual's data \cite{mendes2025synthetic}. Model inversion attacks can reconstruct approximations of training samples from model parameters, potentially exposing sensitive patient information. These attacks prove particularly concerning for rare diseases where small cohorts may enable reidentification even without direct attribute disclosure.

The distinction between fully synthetic and partially synthetic data critically impacts privacy guarantees. Fully synthetic datasets generated without direct reference to real individuals provide stronger privacy protection but may sacrifice analytical validity. Partially synthetic data, which combines real and fabricated elements, maintains higher fidelity to real-world distributions but introduces residual reidentification risks \cite{mendes2025synthetic}. The optimal balance depends on specific use cases, with different applications prioritizing privacy protection versus analytical accuracy differently.

Differential privacy offers theoretical guarantees by adding calibrated noise during training to bound the influence of any individual training sample. However, the privacy-utility trade-off proves challenging for rare diseases where small sample sizes make the impact of individual patients more pronounced. Strong privacy guarantees require noise levels that may substantially degrade model performance, while weak guarantees provide insufficient protection. Choi et al. demonstrated privacy-preserving EHR synthesis without explicit differential privacy mechanisms \cite{choi2017generating}, but the absence of formal privacy bounds limits confidence in the approach's robustness against sophisticated attacks.

Adversarial debiasing and fairness-aware algorithms represent emerging approaches to address representation biases in synthetic data. Mendes et al. noted that synthetic generation can either amplify or mitigate demographic biases present in training data \cite{mendes2025synthetic}. Without explicit fairness constraints, generative models may underrepresent minority populations or rare disease subtypes, perpetuating existing health disparities. Fairness evaluation frameworks for synthetic biomedical data remain nascent, with limited consensus on appropriate metrics or acceptable bias thresholds across different clinical contexts.

\subsection{Computational and Resource Barriers}

The computational demands of deep generative models create substantial accessibility barriers, particularly for smaller research institutions and clinical centers without dedicated machine learning infrastructure. GAN training typically requires GPU clusters, extensive hyperparameter tuning, and expertise in managing training instabilities such as mode collapse. Finetti et al. emphasized that these computational requirements often exceed the resources available to rare disease patient advocacy organizations and smaller medical centers that would most benefit from synthetic data expansion \cite{finetti2025data}.

Training times ranging from hours to days for individual models compound these challenges. Iterative refinement of architectures, hyperparameters, and conditioning mechanisms demands substantial computational budgets. Cloud computing infrastructure provides partial solutions but introduces data governance complexities when protected health information must be processed in third-party environments. The carbon footprint of large-scale model training also raises sustainability concerns as generative model deployment expands.

VAEs offer computational advantages relative to GANs, requiring fewer training iterations and exhibiting more stable convergence \cite{mendes2025synthetic}. However, this efficiency comes at the cost of reduced sample sharpness for imaging applications. Diffusion models, while achieving impressive generation quality, require multiple denoising steps during inference that slow sample generation compared to single-pass GAN or VAE synthesis \cite{zhong2024meddiffusion}. Recent advances in accelerated sampling partially address this limitation but do not eliminate the computational gap.

The expertise requirements extend beyond computational infrastructure to encompass deep learning engineering, domain knowledge integration, and clinical validation skills. Rare disease researchers with medical expertise may lack machine learning training, while machine learning researchers may lack the domain knowledge necessary to assess biological plausibility. This interdisciplinary challenge necessitates collaborative teams that prove difficult to assemble, particularly for ultra-rare conditions affecting dozens of patients globally.

\subsection{Clinical Translation and Regulatory Challenges}

The pathway from research prototypes to regulatory-approved clinical decision support tools remains poorly defined for synthetic data applications. Regulatory agencies including the U.S. Food and Drug Administration (FDA) and European Medicines Agency (EMA) lack established frameworks for evaluating medical devices trained on synthetic data. Key questions remain unresolved: What proportion of training data can be synthetic? How should synthetic data quality be assessed for regulatory purposes? What validation studies demonstrate that models trained on synthetic data generalize safely to real patients?

Clinical integration faces additional barriers beyond regulatory approval. Healthcare providers require interpretability and explainability to trust model recommendations, yet generative models function as black boxes with limited mechanistic transparency. Alsentzer et al. addressed this through attention mechanisms that identify which phenotypic features drive diagnostic predictions \cite{alsentzer2025few}, but such interpretability remains uncommon. Without clear explanations for why a model recommends specific diagnoses or treatments, clinicians may appropriately hesitate to incorporate these tools into practice.

Liability concerns compound adoption challenges. When a diagnostic error occurs with a model trained partially on synthetic data, determining responsibility proves complex. Is the synthetic data generator liable? The model developer? The clinician who applied the recommendation? These legal ambiguities discourage risk-averse healthcare institutions from deploying synthetic data-based systems despite potential clinical benefits.

The generalization of models across different clinical sites, patient populations, and temporal contexts requires extensive validation. Models trained at academic medical centers may fail when deployed at community hospitals with different patient demographics and disease prevalence. Alsentzer et al. demonstrated sustained performance across 12 clinical sites \cite{alsentzer2025few}, but such multi-site validation remains rare. Temporal validation, assessing whether models maintain performance as medical practice evolves and new treatments emerge, receives even less attention despite its critical importance for long-term deployment.

\subsection{Biological Plausibility and Domain Constraint Enforcement}

Generative models trained on purely data-driven objectives lack inherent mechanisms to enforce biological constraints, creating risks of generating physiologically impossible samples. A GAN might produce a medical image with tissue densities outside physically realizable ranges or generate an electronic health record with mutually exclusive diagnosis codes. While such samples may appear statistically consistent with training distributions, they violate domain knowledge that trained medical professionals recognize immediately.

Hybrid approaches that combine generative modeling with rule-based constraints offer partial solutions. Knowledge graphs encoding gene-disease-phenotype relationships, as employed by Alsentzer et al. \cite{alsentzer2025few}, guide generation toward biologically plausible combinations. Ontologies such as the Human Phenotype Ontology or Disease Ontology provide structured representations of valid medical concepts and their relationships, enabling post-generation filtering of invalid samples. However, comprehensively encoding all biological constraints proves infeasible given the complexity of human physiology and the incompleteness of current biomedical knowledge.

The challenge intensifies for novel or poorly characterized rare diseases where biological constraints remain incompletely understood. Generative models might produce samples that appear plausible based on limited existing data yet fail to capture emergent disease characteristics not represented in small training sets. This limitation proves particularly acute for the ultra-rare diseases where generative approaches promise the greatest potential impact but where biological understanding remains most limited.

Validation frameworks must distinguish between statistical fidelity and clinical utility. A synthetic dataset that perfectly reproduces training distribution statistics may nonetheless fail to capture subtle but clinically critical features. Conversely, synthetic data that exhibits some statistical discrepancies from real data might nonetheless provide sufficient signal for training clinically useful predictive models. Establishing which aspects of the real data distribution must be preserved and which can be approximated represents an ongoing challenge requiring close collaboration between machine learning researchers and domain experts.
